\section*{Online supplement}

\subsection*{S1. Data provenance and land base}

The compiled dataset is derived from the public FEMIC TSA24 instance of the UBC FRESH laboratory
\citep{femic_tsa24_2026}, built from the BC Vegetation Resources Inventory (2025), the Forest and
Range Evaluation Program disturbance mapping, and Site Productivity estimates, under the Open
Government Licence -- British Columbia. The compiled dataset carries the inventory tables.
The baseline schedule is the even-flow plan that realises 51.9 million m3 per period. The
standards are re-optimised from that schedule; the enterprise response is described in the main
text.

\begin{table}[h]
    \centering
    \small
    \caption{Compiled estate tables and land-base comparison.}
    \label{tab:supp_estate}
    \begin{tabular}{lr}
        \toprule
        Item & Value \\
        \midrule
        Blocks and fragments                          & 458,708 \\
        Treatments                                    & 48 \\
        Growth curves                                 & 10,326 \\
        Accounts                                      & 99 \\
        Compiled forest area                          & 5,899,679.6 ha \\
        Managed (THLB) area                           & 3,060,473.1 ha \\
        Unmanaged forest area                         & 2,839,206.5 ha \\
        Published THLB, 2017 base case                & 3,070,301 ha \\
        Difference                                    & $-9{,}828$ ha ($-0.32$ percent) \\
        Official allowable annual cut                 & 6,917,231 m3/yr \\
        \bottomrule
    \end{tabular}
\end{table}

The managed area is the sum of block areas whose track carries \texttt{IFM(managed)}; the
unmanaged area is the complement.

\subsection*{S2. Model configuration}

\begin{table}[h]
    \centering
    \small
    \caption{Model configuration.}
    \label{tab:supp_settings}
    \begin{tabular}{ll}
        \toprule
        Setting & Value \\
        \midrule
        Horizon                     & 20 periods of 10 years \\
        Harvest treatment           & clearcut with regeneration \\
        Even-flow target            & within 1 percent of period 1, non-linear, weight 10,000 \\
        Level aspiration            & 70,000,000 m3/period; 51.9 million realised here \\
        Fire                        & recorded regime; composition-dependent severity \\
        \bottomrule
    \end{tabular}
\end{table}

The even-flow rule sets the level; the realised range is 51.89 to 52.06 million m3 per period.

\subsection*{S3. Economic constants}

\begin{table}[h]
    \centering
    \small
    \caption{Economic constants and their sources.}
    \label{tab:supp_constants}
    \begin{tabular}{lll}
        \toprule
        Constant & Value & Source \\
        \midrule
        SPF log price            & 113.96 CAD/m3 & BC Interior Log Market, quarter ending 2026-09-30 \\
        Hem-Bal log price        & 105.98 CAD/m3 & same \\
        Deciduous log price      & 35 to 40, midpoint 37.5 CAD/m3 & 2018 ten-year average \\
        Harvest, tree-length     & 21.38 CAD/m3 & FHOPS machine rates \\
        Haul, 90 km              & 7.85 CAD/m3 & FHOPS highway truck \\
        Silviculture             & 10.02 CAD/m3 & BC Interior costs, North Central zone, 2026-10-01 \\
        Conifer stumpage         & 9.92 CAD/m3 & BC tabular rate, North Central, 2026-11-01 \\
        Deciduous stumpage       & 0.50 CAD/m3 & grades 1 and 2 \\
        Discount rate            & 4 percent real & \citet{bc_mof_silviculture_lesson4} \\
        \bottomrule
    \end{tabular}
\end{table}

Net present value discounts each period's net value at its end. The base rate is the BC Ministry
of Forests 4 percent real rate for public-sector forestry investment analysis; \citet{heaps_pratt_1989} recommend 3 to 5 percent, and the federal social rate for environmental
analysis is 3 percent \citep{treasury_board_2007}. The sensitivity analysis varies the deciduous price over its
documented 35 to 40 CAD/m3 range, the haul distance from 90 to 117 km, and the discount rate over
3, 4, 5, and 8 percent.

\subsection*{S4. Stocking-standard experiments}

The baseline dataset has one treatment, \texttt{CC}, which regenerates the same analysis unit.
The policy dataset adds three prescriptions: \texttt{CC\_D26}, \texttt{CC\_D55}, and
\texttt{CC\_D84}. Each targets a different analysis unit whose managed species-proportion curves
specify the mixture, and the regenerated stand takes the species mix of the prescription. On a conifer block the realised composition and the projected fuel type are:

\begin{table}[h]
    \centering
    \small
    \caption{Stocking standards, realised composition on a conifer block, and the projected fuel
    type and severity multiplier. The multiplier is relative to the initial-landscape mean.}
    \label{tab:supp_standards}
    \begin{tabular}{lrrrr}
        \toprule
        Standard & Target deciduous [\%] & Realised composition & Fuel type & Multiplier \\
        \midrule
        CC  & source & SX 83 / BL 17 & C2/C1 & 1.299 / 0.749 \\
        D26 & 26     & AT 26 / SX 74 & C2/C1 & 1.299 / 0.749 \\
        D55 & 55     & AT 55 / SX 45 & M1    & 1.011 \\
        D84 & 84     & AT 84 / SX 16 & D1    & 0.267 \\
        \bottomrule
    \end{tabular}
\end{table}

A standard is applied only where it increases the deciduous share of the source analysis unit;
aspen blocks stay on \texttt{CC}. The no-fire net present value and volume by standard are in the main text. A planted stand on a pre-harvest fragment with a VRI crown closure below 5 percent
uses a 5 percent floor, because the VRI value is a sparse or non-forest artefact.

The severity multiplier is sensitive to wind and the surface-mortality floor. Across
\SIrange{6.8}{25}{\kilo\meter\per\hour} and a floor of 0.1 to 0.4, the D84 multiplier ranges
0.11 to 0.49, the D55 multiplier 1.01 to 1.04, and the D26 (M2) multiplier 1.03 to 1.24. The D1
multiplier is the model's 0.2 floor, not an estimate of D1 severity. The CVaR effect of the floor
and threshold is reported in the main text; the multiplier range alone does not establish that
the ranking is unchanged.

Over the 720 potential-fire weather states in the 1950--2024 calibration sample, the D1 multiplier
is the 0.267 floor at every state. The C1 multiplier has a median of 1.112 and a 10th-to-90th
percentile range of 0.817 to 1.312; the single-state value of 0.749 is the 10th percentile. The M1
median is 1.033 (range 0.376 to 1.253) and the C2 median is 1.312 (range 1.245 to 1.331). The D84
severity result is therefore weather-robust because it is the floor; the 26 percent standard's
C1/C2 mix is not, and its single-state near-zero aggregate effect is not stable.

\subsection*{S5. Rotation estimate}

For every block with more than one harvest in the schedule, the rotation is the interval between
consecutive harvests. The baseline yields 53,071 intervals, an area-weighted mean of 102.2 years
(median 100, tenth-to-ninetieth percentile 60 to 160). The estimate is conditional on re-entry
within the 200-year horizon and is therefore biased toward productive and accessible stands. It
describes the baseline schedule; the standards are re-optimised and their rotation is an outcome
of the enterprise response. The summaries are in the archived analysis repository; the area-weighted median managed age is
92 years, and 44 percent of the managed area is older than 100 years.

\subsection*{S6. Reproduction}

The derived species-period table, the rotation summaries, the stocking-standard surface, the fuel
projection checks, the fire regime and surface, and the valuation and
rotation code are archived in a public repository with a DOI on acceptance. The sensitivity values
are reproducible from the valuation code and the recorded constants. The policy dataset, the
schedules, the no-fire baseline, and the paired fire ensemble are rebuilt by the scripts in that
repository. The estate model is WS4 \citep{ubc_fresh_ws4_2026}; harvest costing uses FHOPS
\citep{fhops_2026}. The fire calibration and its perimeter set are documented in
\citet{ubc_fresh_fire_calibration_2026}.
